\scanpage{24}
% archive foliation: 23
est nul; or
\[
q^*(e_{i,m})=\unsure{m|qm}e_i,\qquad
p_*(e_{i,\unsure{m|n}})=\unsure{n|pn}e_i
\]
où $e_i$ définit l'extension $\Phi^i$.
Donc la condition est que
\[
(qm-pn)e_i=0
\]
On fait l'hypothèse
\[
\text{(c)}\qquad
\boxed{e_i\text{ est une base de }\operatorname{Ext}^1_{C_0}(\Psi^i,\Psi^{i-1})}
\]
de sorte que la condition devient
\[
qm=pn
\]
i.e.\quad
\[
\frac q p=\frac{\unsure{m|n}}{\unsure{n|m}}
=\frac{\unsure{m'|n'}}{\unsure{n'|m'}}
\]
si $m'=\dfrac{m}{(m,n)}$, $n'=\dfrac{n}{(m,n)}$, donc
on aura
\[
p=\deleted{\unsure{pq}}\,t m',\quad q=t n'\qquad(t\in k).
\]
Or pour $p,q$ correspondant ainsi
à un $t$ (évidemment uniquement
déterminé par $(p,q)$) l'hom
$\alpha$ qui leur correspond est déterminé
mod un $\operatorname{Hom}(\Psi^i,\Psi^{i-1})=0$,
donc est unique. On trouve donc
\[
\operatorname{Hom}(\Phi(i,n),\Phi(i,m))\simeq k.
\]
